% Some LCHS basics if needed
\section{Appendix}
\subsection{Complementary (Homogeneous) Solution  $\Tcal e^{- \int^t_0 A(s) \dd s}u_0$}

\begin{remark}
With a quadrature rule, the goal is to find the approximation
\begin{align}
    \int^b_a g(x) \approx \sum_{j=0}^{m-1} w_j g(x_j).
\end{align}
For an integral over $[-1,1]$ with Gaussian quadrature with $m$ nodes, set node $x_j$ to be the root of Legendre polynomial of degree $m$, $P_m(x)$ with weight $w_j = \frac{2}{(1-x^2_j)((P'_m(x_j))^2}$. For a longer interval $[a,b]$, the composite rule is applied by dividing the interval $[a,b]$ into $K$ segments, i.e., $[a+kh, a+(k+1)h]$ where $h = \frac{b-a}{K}$ is the step size. Gaussian quadrature is applied on each shorter interval and it is constructed from that on $[-1,1]$ by a change of variable. That is, the weight $w_j \rightarrow \frac{h}{2}w_j$ and node $x_j \rightarrow \frac{h}{2}x_j + (a+kh) + \frac{1}{2}h$.
\end{remark}


Consider the ODE solution over a truncated finite interval $[-K,K]$ using the composite Gaussian quadrature to discretize the interval. 

The first step is to define
\begin{align}
    g(k) = \frac{1}{C_\beta (1-ik)e^{(1+ik)^\beta}}, \quad U(T,k) = \Tcal e^{-i \int^T_0 (kL(s)+H(s)) \dd s}.
\end{align}
Then,
\begin{align}
    \Tcal e^{-\int^T_0 A(s) \dd s} &= \int_\Rbb g(k) U(T,k) \dd k   && \text{(by LCHS)}    \nonumber \\
    &\approx \int^K_{-K} g(k) U(T,k) \dd k      && \text{(truncation)} \nonumber \\
    &= \sum_{m = -K/h_1}^{K/h_1-1} \int^{(m+1)h_1}_{mh_1} g(k)U(T,k) \dd k    && \text{(divide to shorter intervals)}\nonumber \\
    &\approx \sum_{m=-K/h_1}^{K/h_1 -1}\sum_{q=0}^{Q-1} c_{q,m}U(T,k_{q,m})   && \text{(Gaussian quadrature in each interval)} \label{eq:homo-gauss1}\\
    &= \sum_{j=0}^{M-1} c_j U(T,k_j)  && \text{(simplified notation)} 
\end{align}
where $h_1$ is the step size in the composite Gaussian quadrature, $Q$ is the number of nodes in each interval $[mh_1, (m+1)h_1]$, $k_{q,m}$ is a Gaussian node, $c_{q,m} = w_qg(k_{q,m})$, $w_q$ is a Gaussian weight (which does not depend on the choice of $m$), and $M$ is the overall number of the nodes (or unitaries). 

% \begin{align}
%     \Tcal e^{-\int^T_0 A(s) \dd s} = \int_\Rbb g(k) U(T,k) \dd k  \approx \int^K_{-K} g(k) U(T,k) \dd k     = \sum_{j=0}^{M-1} c_j U(T,k_j) 
% \end{align}

\begin{assumption}
    Assume $K$ and $h_1$ are chosen so that $K/h_1$ is an integer.
\end{assumption}

To bound the truncation error by $\epsilon > 0$, it suffices to choose 
\begin{align*}
    K = \Ocal \left(\left(\log(1/\epsilon)\right)^{1/\beta}\right)
\end{align*}
To bound the quadrature error by $\epsilon > 0$, it suffices to choose
\begin{align}
    h_1 = \frac{1}{eT \max_t \|L(t)\|}, \ Q = \left\lceil \frac{1}{\log 4}  \log\left( \frac{8}{3C_\beta} \frac{K}{\epsilon} \right)\right\rceil = \Ocal\left(\log \frac{1}{\epsilon}\right), \label{eq:homo-hq}
\end{align}
and the overall number of unitaries in the summation formula is
\begin{align}
    M = \frac{2KQ}{h_1} = \Ocal\left( T \max_t \|L(t)\|\left(\log \frac{1}{\epsilon}\right)^{1+1/\beta} \right)
\end{align}




For the specific value of $K$, note that in Eq. (176) in~\cite{LCHS2}, we have 
\begin{align}
\left\|\int_\Rbb g(k) U(T,k) \dd k  -\int^K_{-K} g(k) U(T,k) \dd k  \right\| &\leq \frac{2^{B+1}B!}{C_\beta (\cos(\beta \pi /2))^B} \frac{1}{K} e^{-\frac{1}{2}K^\beta \cos(\beta \pi /2)} \leq \epsilon \nonumber \\
\ln \left( \frac{2^{B+1}B!}{C_\beta (\cos(\beta \pi /2))^B}\right)  + \ln(1/ \epsilon)  &\leq \ln(K) + \frac{1}{2}K^\beta \cos(\beta \pi /2) \\
2\cdot \frac{\ln \left( \frac{2^{B+1}B!}{C_\beta (\cos(\beta \pi /2))^B}\right)  + \ln(1/ \epsilon) - \ln(K)}{ \cos(\beta \pi /2)} &\leq K^\beta
\end{align}
where $B = \lceil 1/\beta \rceil$








%%%%%%%%%%%%%%%%%%%

\subsection{Inhomogeneous Term $\int^t_0 \Tcal e^{-\int^t_s A(s')\dd s'}b(s) \dd s$}


The discretization of the inhomogeneous term is
{\footnotesize
\begin{align}
    &\quad \int^T_0 \Tcal e^{-\int^T_s A(s')\dd s'}b(s) \dd s  \nonumber \\
    &= \int^T_0 \int_\Rbb \frac{f(k)}{1-ik} U(T,s,k) b(s) \dd k \dd s   && \text{(by LCHS Eq.~\eqref{eq:LCHS-sol-homo})}    \nonumber \\
    &\approx \int^T_0  \sum_{m_1 = -K/h_1}^{K/h_1 - 1} \sum_{q_1 = 0}^{Q_1 - 1} c_{q_1, m_1} U(T, s, k_{q_1, m_1}) b(s) \dd s        &&\text{(by 
 Eq.~\eqref{eq:homo-gauss1})}   \nonumber \\
    &\approx \sum_{m_2 = 0}^{T/h_2 - 1} \sum_{q_2 = 0}^{Q_2 - 1}    \sum_{m_1 = -K/h_1}^{K/h_1 - 1} \sum_{1_1 = 0}^{Q_1 - 1} c'_{q_2, m_2} c_{q_1, m_1} U(T, s_{q_2, m_2}, k_{q_1, m_1} ) \ket{b(s_{q_2, m_2})}        &&\text{(composite Gaussian quadrature) } \\
    &= \sum_{j'=0}^{M'-1} \sum_{j = 0}^{M-1} c'_{j'} c_j U(T, s_{j'}, k_j)\ket{b(s_{j'})} && \text{(simplified notation)} \label{eq:inhomo}
\end{align}
}
where $ U(T,s,k) = \Tcal e^{-i\int^T_s (kL(s') + H(s') \dd s' }$, $h_2$ is the step size for the discretization of the time integral, $Q_2$ is the number of Gaussian quadrature nodes, $s_{q_2, m_2}$ is a node, $c'_{q_2,m_2} = w'_{q_2}\|b(s_{q_2, m_2})\|$, the $w'_{q_2}$ is a Gaussian weight. The $\ket{b}$ implies that it is the normalized vector $b/\|b\|$.

Let 
\begin{itemize}
    \item $\Lambda = \sup_{p \geq 0, t\in [0,T]} \|A^{(p)}\|^{1/(p+1)}$ 
    \item $\Xi = \sup_{p \geq 0, t\in [0,T]} \|b^{(p)}\|^{1/(p+1)}$, where $(p)$ refers to the $p^{th}$-order time derivative
    \item $\|b\|_{L^1} = \int^T_0 \|b(s)\| \dd s$
\end{itemize}
To bound the approximation error of Eq.~\eqref{eq:inhomo} by $\epsilon > 0$, it suffices to choose
\begin{align}
    K = \Ocal\left( \left( \log\left( 1+ \frac{\|b\|_{L^1}}{\epsilon}\right) \right)^{1/\beta} \right), \   h_1 = \frac{1}{eT\max_t \|L(t)\|}, h_2 = \frac{1}{eK(\Lambda+\Xi)}, \label{eq:inhomo-K} \\
    Q_1 = \Ocal\left( \log\left( 1+\frac{\|b\|_{L^1}}{\epsilon} \right) \right), Q_2 = \Ocal\left( \log\left( \frac{T(\Lambda+\Xi)}{\epsilon} \right) \right)
\end{align}
and
\begin{align*}
    M' = \frac{TQ_2}{h_2} = \Ocal\left( TK(\Lambda+\Xi)\log\left( \frac{T(\Lambda+\Xi)}{\epsilon} \right)  \right) = \tilde{\Ocal}\left( T(\Lambda+\Xi) \left( \log\left( 1+ \frac{\|b\|_{L^1}}{\epsilon}\right) \right)^{1/\beta} \log\left(1/\epsilon\right)\right)
\end{align*}

As indicated in Eq. (234) and (235) in~\cite{LCHS2}, we have
\begin{align}
    \|c\|_1 \frac{eT(\Lambda+\Xi)}{4^{Q_2} } &\leq \frac{\epsilon}{2} \nonumber \\
    \frac{1}{ \ln(4)} \ln \left( \|c\|_1 \frac{2eT(\Lambda+\Xi)}{\epsilon} \right) &\leq Q_2,
\end{align}
where $ \|c\|_1 = \Ocal(1)$. For computing $Q_1$, just use the same formula in Eq.~\eqref{eq:homo-hq} but with the $K$ in Eq.~\eqref{eq:inhomo-K}.


\subsection{Summary}
With LCHS, Eq.~\eqref{eq:exp-sol} becomes
\begin{align}
    u(T) \approx \sum_{j=0}^{M-1} c_j U(T,k_j) \ket{u_0}  + \sum_{j'=0}^{M'-1} \sum_{j = 0}^{M-1} c'_{j'} c_j U(T, s_{j'}, k_j)\ket{b(s_{j'})}.
\end{align}
where $ U(T,s,k) = \Tcal e^{-i\int^T_s (kL(s') + H(s')) \dd s' }$ and $U(T,k) = U(T,0,k)$. So, we need oracles circuit to implement the Hamiltonian simulation operator $U(T,s,k)$ and  oracles to encode the coefficients $c$ and $c'$ into amplitudes. Then, use LCU to obtain $u(T)$.




\subsection{Analytical Values of the Expansion Coefficients $C_n$ at limit values of $\beta$}
\label{ss:analytic}
The coefficients are defined by the projection of the signal $g(x)$ onto the Hermite basis:

\begin{equation}
    C_n(\beta) = \mathcal{K}_n \int_{-\infty}^{\infty} \frac{H_n(x/\sigma') e^{-\gamma x^2}}{2\pi e^{-2^\beta} e^{(1+ix)^\beta} (1-ix)} \, \dd x
\end{equation}

where $\mathcal{K}_n = \sqrt{\frac{\sigma}{\sigma'}} \frac{1}{\sqrt{2^n n!}}$ and $\gamma = \sigma'^{-2} - \sigma^{-2}$. We denote the coefficient functions at the limit points as $C_n^{(0)}$ and $C_n^{(1)}$, corresponding to the limits $\beta=0$ and $\beta=1$, respectively.

We employ an explicit polynomial decomposition of the Hermite-Lorentzian kernel. By defining a residual quotient polynomial $Q_{n-1}(x) = [H_n(x/\sigma') - H_n(-i/\sigma')] / (x + i)$, we isolate the singular contribution of the complex pole at $x = -i$.

At the limit $\beta = 0$, the denominator simplifies to a standard complex Lorentzian. The exact closed-form expression is:
\begin{equation}
    C_n^{(0)} = \mathcal{K}_n \left[ \frac{1}{2} H_n\left(-\frac{i}{\sigma'}\right) e^\gamma \text{erfc}(\sqrt{\gamma}) + \frac{i}{2\pi} \int_{-\infty}^{\infty} Q_{n-1}(x) e^{-\gamma x^2} \, \dd x \right]
\end{equation}

For $\beta = 1$, the kernel incorporates a phase shift $e^{-ix}$ and a scaling factor $e^1$. Completing the square in the exponent for the shifted Gaussian ($-\gamma x^2 - ix$) leads to:
\begin{equation}
    C_n^{(1)} = \mathcal{K}_n \left[ \frac{1}{2} H_n\left(-\frac{i}{\sigma'}\right) e^\gamma \text{erfc}\left(\sqrt{\gamma} - \frac{1}{2\sqrt{\gamma}}\right) + \frac{ie}{2\pi} \int_{-\infty}^{\infty} Q_{n-1}(x) e^{-\gamma x^2 - ix} \, \dd x \right]
\end{equation}
The transition from $\beta=0$ to $\beta=1$ effectively shifts the argument of the error function by $\Delta = (2\sqrt{\gamma})^{-1}$.




\subsubsection{Numerical Validation}
The analytical results were compared against direct numerical quadrature using a double-exponential integration scheme with a working precision of 30 decimal places. For the parameter set $\sigma' = 1.2$ and $\sigma = 2.0$ ($\gamma = 4/9$), the agreement is presented in the tables below.




\begin{table}[h!]
\centering
\caption{Comparison of Numerical and Analytical Coefficients for the Lower Bound ($\beta = 0, \gamma = 4/9$)}
\label{tab:lower_bounds}
\begin{tabular}{rll}
\hline
$n$ & Numerical Result & Analytical (Explicit) \\
\hline
0 & 0.3481065885 & 0.3481065885 \\
1 & $0.2335427879i$ & $0.2335427879i$ \\
2 & $-0.0515295394$ & $-0.0515295394$ \\
3 & $-0.0077856810i$ & $-0.0077856810i$ \\
\hline
\end{tabular}
\end{table}





\begin{table}[h!]
\centering
\caption{Comparison of Numerical and Analytical Coefficients for the Upper Bound ($\beta = 1, \gamma = 4/9$)}
\label{tab:upper_bounds}
\begin{tabular}{rll}
\hline
$n$ & Numerical Result & Analytical (Explicit) \\
\hline
0 & 1.1011786413 & 1.1011786413 \\
1 & $-0.3006294949i$ & $-0.3006294949i$ \\
2 & $-0.0943736080$ & $-0.0943736080$ \\
3 & $-0.1769037943i$ & $-0.1769037943i$ \\
\hline
\end{tabular}
\end{table}



As demonstrated in Tables~\ref{tab:lower_bounds} and~\ref{tab:upper_bounds}, the explicit decomposition method maintains a relative error below $10^{-25}$.


\subsection{Translation to the $\hbar=1$ Backend Convention}
\cb{The main text uses the $\hbar=2$ convention. For numerical backend convenience, we record the corresponding $\hbar=1$ formulas used by the QuTiP and bosonic simulation codes.}

\cb{If the paper convention is written with}
{\color{blue}
\[
\hat x_{(2)} = a+a^\dagger,
\]
}
\cb{then the backend convention is}
{\color{blue}
\[
\hat x_{(1)}=\frac{a+a^\dagger}{\sqrt{2}},
\qquad
\hat x_{(2)}=\sqrt{2}\,\hat x_{(1)}.
\]
}

\cb{Accordingly, the coefficient-exponent parameter appearing in the prepared-state overlap changes from}
{\color{blue}
\[
\gamma_{(2)} = e^{-2r'}-e^{-2r}
\]
}
\cb{to}
{\color{blue}
\[
\gamma_{(1)} = \frac{1}{2}\left(e^{-2r'}-e^{-2r}\right)=\frac{\gamma_{(2)}}{2}.
\]
}

\cb{Thus the backend coefficient integral is}
{\color{blue}
\begin{equation}
C_n^{(\hbar=1)} \propto \int_{-\infty}^{\infty}
H_n\!\left(\frac{k}{\sigma'}\right)\,
g_\beta(k)\,
e^{-\gamma_{(1)} k^2}\,dk,
\qquad
\gamma_{(1)}=\frac12\left(e^{-2r'}-e^{-2r}\right).
\end{equation}
}

\cb{Likewise, a hybrid interaction written in the paper convention as}
{\color{blue}
\[
e^{-i\lambda\,\hat x_{(2)}\otimes P}
\]
}
\cb{is implemented in an $\hbar=1$ backend as}
{\color{blue}
\[
e^{-i(\sqrt{2}\lambda)\,\hat x_{(1)}\otimes P},
\]
}
\cb{so the coupling coefficient supplied to the backend must be rescaled by a factor of $\sqrt{2}$.}


\section{Fourier-Laplace Transform}
\cr{The Fourier-Laplace section below is not yet theoretically consistent: the kernel notation mixes $k$ and $z$, the normalization is malformed, and the two-oscillator theorem does not yet match the operator structure being claimed.}
\cb{For the current working-paper version, this section should be treated as exploratory only. A clean revision should use one integration variable, fix the kernel normalization, and derive the two-oscillator sandwich directly from the target Laplace-transform representation before stating a theorem.}

Ref.~\cite{an2024laplacetransformbasedquantum} proposed an LCHS-based Laplace transform QEVT. The theorem is as follows:

Let 
\[
A = L + iH, \quad \text{where} \quad 
L = \frac{A + A^{\dagger}}{2} \succeq 0
\quad \text{and} \quad 
H = \frac{A - A^{\dagger}}{2i}
\]
are both Hermitian matrices. 

For a function 
\[
h(z) = \int_{0}^{\infty} g(t)e^{-zt}\,dt
\]
that can be expressed by a Laplace transformation with \(g(t)\) being its inverse Laplace transform, we have
\[
h(A) = \int_{0}^{\infty} \int_{\mathbb{R}} 
\frac{f(k)g(t)}{1 - ik}\, e^{-it(kL+H)} \, dk \, dt.
\tag{5}
\]

Here
{\color{red}
\[
f(k) = \frac{1}{2\pi e^{-2\beta} e^{(1+iz)\beta}}, 
\qquad \beta \in (0,1).
\]
}
\cb{If this appendix keeps the same near-optimal kernel convention as the main text, it should instead use}
{\color{blue}
\[
f_\beta(k)=\frac{e^{2^\beta}}{2\pi e^{(1+ik)^\beta}},
\qquad
g_\beta(k)=\frac{f_\beta(k)}{1-ik},
\qquad
\beta \in (0,1),
\]
}
\cb{with the same branch choice and normalization used throughout the rest of the paper.}

We propose the CVDV LCHS based Laplace transform theorem as follows:
{\color{red}
\begin{theorem}
\begin{align}
    \ket{u(t)} &= h(A) \ket{u(0)}= \bra{\Phi_{\rm osc_1}}\braket{\Phi_{\rm osc_2}|e^{-i(\hat x_1\otimes \hat x_2\otimes L+\hat{x}_1\otimes \mathbf{I}\otimes H)}|\Psi_{\rm osc_2}}\ket{\Psi_{\rm osc_1}} \\
    &= (\bra{\phi}_{\rm osc_1}\otimes\mathbf{I_{\rm osc_2}}\otimes\mathbf{I}_q) (\mathbf{I_{\rm osc_1}}\otimes\bra{\phi}_{\rm osc_2}\otimes\mathbf{I}_q) e^{-it ( \hat{x} \otimes L+ \mathbf{I}_{\rm osc} \otimes H} (\ket{\psi}_{\rm osc_1}\otimes\ket{\psi}_{\rm osc_2}\otimes\ket{u_0}_q)\\
    &= \left[\int_{\mathbb{R}} \int_{\mathbb{R}} \dd x_1 \dd x_2 \, \phi_1^*(x_1) \psi_1(x_1)   \, \phi_2^*(x_2) \psi_2(x_2)   e^{-it ( xL + H)}\right]\ket{u_0}_q.
\end{align}
\end{theorem}
}
\cb{(...) Replace the theorem above with a fully derived two-oscillator statement after the target Laplace-transform representation is fixed. The revised theorem should specify the joint unitary, both post-selection kernels, the resulting double-integral weight, and the exact conditions under which the construction reproduces $h(A)$.}
